\documentclass[11pt]{article}

\usepackage[T1]{fontenc}
\usepackage{amsmath,amssymb,amsthm,mathtools}
\usepackage[margin=1in]{geometry}
\usepackage{hyperref}

\newtheorem{theorem}{Theorem}
\newtheorem{lemma}[theorem]{Lemma}
\newcommand{\F}{\mathbb{F}}
\newcommand{\Tr}{\operatorname{Tr}}
\newcommand{\MCM}{\operatorname{MCM}}
\newcommand{\Kas}{\operatorname{Kas}}

\title{The Generalized MCM Polynomial Is a Permutation Polynomial}
\author{}
\date{}

\begin{document}
\maketitle

\section{Statement}

Let $F=\F_{2^n}$, and let $0<k<n$.  For $b\in\F_2$, define
\[
  \psi_b(x)=\sum_{i=0}^{k-1}x^{2^i}+b\Tr(x)
\]
and define the generalized Muller--Cohen--Matthews polynomial by
\[
  \MCM_b(X)=
  \left(\sum_{i=0}^{k-1}X^{2^i}+b\Tr(X)\right)^{2^k+1}
  X^{(2^n-1)-2^k}.
\]
Here $\Tr(x)=\sum_{i=0}^{n-1}x^{2^i}$ is the absolute trace of $F$ over
$\F_2$.  This is the polynomial called \texttt{genMCMpol} (with arguments
$n$, $k$, and $b$) in the Lean development.

\begin{theorem}
\label{thm:mcm-permutation}
Suppose that $\gcd(k,n)=1$ and that $k$ is odd.  Then $\MCM_0$ induces a
permutation of $F$.
\end{theorem}

The proof is the $b=0$ specialization of the transfer theorem formalized as
\verb|dobbertin_theorem4|.  We state precisely the part of that theorem used
below and then verify its arithmetic hypotheses.

\section{The transfer theorem}

For positive integers $k,n$ with $\gcd(k,n)=1$, choose integers
$k^{-1}$ and $n'$ satisfying
\begin{equation}
\label{eq:bezout-positive}
  0<k^{-1}<n,
  \qquad
  k^{-1}k=n'n+1.
\end{equation}
For $a\in\F_2$, put
\[
  \phi_a(x)=\sum_{i=0}^{k^{-1}-1}x^{2^{ik}}+a\Tr(x).
\]

\begin{lemma}[Kasami parity criterion]
\label{lem:kasami-parity}
Assume \eqref{eq:bezout-positive} and
\begin{equation}
\label{eq:kasami-parity}
  k^{-1}+an\equiv1\pmod2.
\end{equation}
Then the generalized Kasami polynomial $\Kas_a$ is a permutation polynomial
of $F$.
\end{lemma}

\begin{proof}[Proof idea]
The formal lemma \verb|parity_implies_permutation| proves that the evaluation
map of $\Kas_a$ is injective and then uses finiteness of $F$ to obtain
surjectivity.  Its injectivity argument translates equality of two values of
$\Kas_a$ into two instances of the associated Kasami equation
\[
  c z^{2^k+1}=R_a(z),
\]
where $R_a$ is the equation's right-hand side.  The zero and nonzero input
cases are treated separately so that the required nonzero admissibility
condition is retained when $c\ne0$.

The key uniqueness lemma splits into three cases.  If $c=0$, a separate
zero-right-hand-side result gives the unique solution.  If $c\ne0$ and
$c=\gamma^{2^k+1}+\gamma$ for some $\gamma$, the odd parity
\eqref{eq:kasami-parity} excludes a solution in the zero fiber of the
auxiliary gamma polynomial.  The remaining fiber is controlled by the
Delta and Ell identities; two distinct solutions would force both a parity
identity $e_0+e_1=1$ and the vanishing of the corresponding $e$-terms, a
contradiction.  Finally, when no such $\gamma$ exists, uniqueness follows
from the linearized equation.  Thus every admissible Kasami-equation fiber
has at most one solution, proving injectivity of $\Kas_a$.
\end{proof}

\begin{theorem}[Dobbertin transfer theorem]
\label{thm:transfer}
Assume \eqref{eq:bezout-positive}, and suppose that $a,b\in\F_2$ satisfy
\begin{align}
\label{eq:q-parity}
  k^{-1}+an &\equiv 1 \pmod 2,\\
\label{eq:beta-parity}
  b &\equiv n'+ak \pmod 2.
\end{align}
Then $\psi_b$ and $\phi_a$ are inverse permutations of $F$.  Moreover, if
$\Kas_a$ denotes the associated generalized Kasami polynomial, then
\begin{align*}
  \Kas_a\circ\psi_b &= \iota\circ\MCM_b,\\
  \MCM_b\circ\phi_a &= \iota\circ\Kas_a,
\end{align*}
where $\iota(x)=x^{-1}$, with $\iota(0)=0$.  If the parity condition
\eqref{eq:q-parity} holds, then Lemma \ref{lem:kasami-parity} shows that
$\Kas_a$ is a permutation polynomial;
consequently, $\MCM_b$ is a permutation polynomial.
\end{theorem}

For context, the formal proof of this transfer theorem has two main inputs.
First, Frobenius periodicity $x^{2^n}=x$ and trace invariance show that
$\psi_b$ and $\phi_a$ are mutual inverses under
\eqref{eq:q-parity}--\eqref{eq:beta-parity}.  Second, the parity condition
makes the generalized Kasami polynomial a permutation by the fiber-uniqueness
argument in Lemma \ref{lem:kasami-parity}.  The second displayed functional
identity then makes $\MCM_b\circ\phi_a$ a composition of permutations.
Since $\phi_a$ is itself a permutation, $\MCM_b$ is a permutation as well.

\section{Proof of Theorem \ref{thm:mcm-permutation}}

Because $\gcd(k,n)=1$, $k$ has a multiplicative inverse modulo $n$.  Choose
$k^{-1}$ with
\[
  0<k^{-1}<n,
  \qquad
  k^{-1}k\equiv1\pmod n.
\]
Set
\[
  n'=\frac{k^{-1}k}{n}.
\]
The modular-inverse congruence says that $k^{-1}k$ has remainder $1$ upon
division by $n$, so the division algorithm gives
\begin{equation}
\label{eq:product}
  k^{-1}k=n'n+1.
\end{equation}
Thus the first set of hypotheses of Theorem \ref{thm:transfer} is satisfied.

Choose $a\in\F_2$ to be the parity of $n'$; equivalently,
\begin{equation}
\label{eq:a-parity}
  a\equiv n'\pmod2.
\end{equation}
Since $k$ is odd,
\begin{equation}
\label{eq:k-odd}
  k\equiv1\pmod2.
\end{equation}
We verify the two parity relations required by the transfer theorem.

Multiplying \eqref{eq:k-odd} by $k^{-1}$ and using \eqref{eq:product}
yields
\[
  k^{-1}\equiv k^{-1}k\equiv n'n+1\pmod2.
\]
Together with \eqref{eq:a-parity}, this gives
\[
  k^{-1}+an
  \equiv (n'n+1)+n'n
  =2n'n+1
  \equiv1\pmod2.
\]
Hence \eqref{eq:q-parity} holds.

For the second relation, \eqref{eq:a-parity} and \eqref{eq:k-odd} imply
\[
  ak\equiv n'k\equiv n'\pmod2.
\]
Therefore
\[
  n'+ak\equiv n'+n'\equiv0\pmod2.
\]
This is exactly \eqref{eq:beta-parity} with $b=0$.

All hypotheses of Theorem \ref{thm:transfer} now hold for the data
$(n,k,k^{-1},n',a,b=0)$.  Its conclusion states that $\MCM_0$ is a
permutation polynomial of $F$, which proves Theorem \ref{thm:mcm-permutation}.
\qed

\section{Formalization Correspondence}

The construction of $k^{-1}$ is
\verb|Nat.exists_mul_mod_eq_one_of_coprime|.  The quotient $n'$ is
\verb|nPrime|, equation \eqref{eq:product} is \verb|hprod|, and the element
$a$ is the Lean value \verb|a : Fin 2|.  The two parity calculations are
\verb|hqpar| and \verb|hbeta|.  Lemma \ref{lem:kasami-parity} corresponds to
\verb|parity_implies_permutation|: it uses \verb|kasamiEquation_unique| to
prove injectivity and \verb|Finite.injective_iff_surjective| to conclude
bijectivity.  Finally, the formal theorem invokes
\verb|dobbertin_theorem4| with $b=0$ and selects its final conclusion,
\verb|IsPermutationPolynomial (genMCMpol n k 0)|.

\end{document}
